%! Function Model Selection and Assumption Articulation — Unit 1, Lesson 1.13 — Group Activity
\documentclass[10pt]{article}
\usepackage{precalculus-article}
\usepackage{precalculus-boxes}
\newcommand{\TallMath}[1]{$\displaystyle #1\rule[-1.4em]{0pt}{3.2em}$}

\begin{document}

\pageheader{Unit 1, Lesson 1.13}{Group Activity}
\namepartnerperiod

\vspace{0.08in}

\begin{scenariobox}[Data Dispatch Agency]{sky}
\small
Your team works at the Data Dispatch Agency. Four requests just arrived, each describing a
real-world situation that needs a function model. For each dispatch, your team must (1) select
the best model type, and (2) state the key assumptions and domain restriction. Work through
the tier that matches your group's assignment.
\end{scenariobox}

\vspace{0.08in}

\textbf{The four dispatches:}

\begin{itemize}[itemsep=3pt]
  \item \textbf{Dispatch 1.} A city tracks its population each year. The population increases by
        approximately 850 people per year for the past decade.
  \item \textbf{Dispatch 2.} A farmer wants to fence a rectangular pen using 200 m of fencing.
        She records the area for different choices of one side length.
  \item \textbf{Dispatch 3.} A warehouse stores boxes stacked in cubic arrangements.
        The total number of boxes for a stack of side length $n$ units is recorded.
  \item \textbf{Dispatch 4.} A ride-share app charges \$2.50 per mile for the first 5 miles,
        then \$1.75 per mile for each mile after that.
\end{itemize}

\vspace{0.08in}

\begin{tcolorbox}[colback=white, colframe=black!40,
    title=\textbf{Tier R --- Remediate},
    fonttitle=\bfseries, arc=2mm, left=3mm, right=3mm, top=2mm, bottom=2mm]

For each dispatch, \textbf{circle the best model type} from the three choices given, then
complete the assumption sentence.

\renewcommand{\arraystretch}{1.8}
\begin{tabularx}{\linewidth}{@{} c X p{5.5cm} @{}}
\textbf{Dispatch} & \textbf{Best model (circle one)} & \textbf{Complete the assumption} \\
\midrule
1 & Linear \quad Quadratic \quad Cubic & ``This model assumes the population grows by a \blank{2.5cm} amount each year.'' \\
2 & Linear \quad Quadratic \quad Cubic & ``This model assumes the pen is a \blank{2cm} and the perimeter is exactly \blank{1.5cm} m.'' \\
3 & Linear \quad Quadratic \quad Cubic & ``This model assumes the stacks form perfect \blank{2cm} shapes.'' \\
4 & Linear \quad Quadratic \quad Piecewise & ``This model assumes the rate changes at exactly \blank{1.5cm} miles.'' \\
\end{tabularx}
\end{tcolorbox}

\vspace{0.08in}

\begin{tcolorbox}[colback=white, colframe=black!40,
    title=\textbf{Tier A --- Apply},
    fonttitle=\bfseries, arc=2mm, left=3mm, right=3mm, top=2mm, bottom=2mm]

For each dispatch, compute differences from the table provided, determine the degree, name
the model type, and state one assumption with domain restriction.

\medskip
\textbf{Dispatch 1 table:}
\quad
\renewcommand{\arraystretch}{1.3}
\begin{tabular}{|c|c|c|}
\hline
Year ($x$) & Population ($y$) & $\Delta y$ \\
\hline
0 & 12{,}000 & --- \\
\hline
1 & 12{,}850 & \blank{1.8cm} \\
\hline
2 & 13{,}700 & \blank{1.8cm} \\
\hline
3 & 14{,}550 & \blank{1.8cm} \\
\hline
\end{tabular}

\medskip
Model type: \blank{3cm} \quad Degree: \blank{1cm}

\medskip
Assumption: \writeline

Domain restriction (give an inequality): \blank{5cm}

\bigskip
\textbf{Dispatch 2 table:} (side length $\ell$ in meters, area $A$ in m$^2$)
\quad
\renewcommand{\arraystretch}{1.3}
\begin{tabular}{|c|c|c|c|}
\hline
$\ell$ & $A$ & $\Delta A$ & $\Delta^2 A$ \\
\hline
10 & 900  & ---         & ---         \\
\hline
20 & 1600 & \blank{1.5cm} & ---      \\
\hline
30 & 2100 & \blank{1.5cm} & \blank{1.5cm} \\
\hline
40 & 2400 & \blank{1.5cm} & \blank{1.5cm} \\
\hline
50 & 2500 & \blank{1.5cm} & \blank{1.5cm} \\
\hline
\end{tabular}

\medskip
Model type: \blank{3cm} \quad Degree: \blank{1cm}

\medskip
Domain restriction: \blank{6cm}
\end{tcolorbox}

\vspace{0.08in}

\begin{tcolorbox}[colback=white, colframe=black!40,
    title=\textbf{Tier E --- Extend},
    fonttitle=\bfseries, arc=2mm, left=3mm, right=3mm, top=2mm, bottom=2mm]

\begin{enumerate}[itemsep=10pt]

  \item \textbf{Dispatch 4 --- Piecewise model.} Write a two-piece description of the ride-share
        cost function $C(m)$ where $m$ is the number of miles driven.

        \medskip
        Piece 1 (applies when \blank{3cm}): $C(m) = $ \blank{5cm}

        \vspace{0.05in}
        Function type for Piece 1: \blank{3.5cm}

        \medskip
        Piece 2 (applies when \blank{3cm}): $C(m) = $ \blank{5cm}

        \vspace{0.05in}
        Function type for Piece 2: \blank{3.5cm}

  \item State \textbf{two assumptions} embedded in the Dispatch 4 model and explain
        why each assumption is needed for the model to be valid.

        \medskip
        Assumption 1: \writeline
        Why it matters: \writeline

        \medskip
        Assumption 2: \writeline
        Why it matters: \writeline

  \item For Dispatch 3, a student claims a cubic model is appropriate because ``volume is
        three-dimensional.'' Explain in one or two sentences whether this reasoning is
        mathematically sound, referencing differences.

        \medskip
        \writelines{2}

\end{enumerate}
\end{tcolorbox}

\end{document}
